\noindent
This paper addresses the question prior to that of its predecessor. The prior paper examined whether a generative cycle is good and whether it is sustainable; the present paper examines the antecedent question of how relation and value are constituted at all, how a value not yet named comes to be, and how the language by which a bond subsists is made. The method proceeds in order of grounding. It begins in phenomenology, with the datum, prior to theory, in which two persons co-apprehend a value generated between them and not yet symbolised. From the description of that datum it derives an ontology (relation as the being of a generative system, grammar as the manifestation of the system's dynamics, value as its phenomenon) in which the phenomenon reflexively reconfigures the dynamics that produce it, and the subject is generated within the cycle it generates. The formal epistemologies are admitted only thereafter: a quantum-structural relational grammar, employed strictly as an epistemic model and not as a physical hypothesis, in which co-apprehension is a participatory measurement that generates rather than registers its value; and, symmetrically, the political-economic thesis that the circulation of value reconfigures the subject that creates it, alienation and self-realisation being the negative and positive phase of one reflexive cycle, distinguished by the sign of an independently definable reproduction surplus. These are distinct symmetry-broken phases of one abstract structure, which no framework possesses externally; their plurality is established as the method dialectical materialism requires of a totality not yet complete, rather than as eclecticism. The paper concludes in a practice and an ethics. The practice is not the construction of a grammar but the cultivation of a field within which a grammar admits of continual joint reconfiguration; and the ethics turns on the result that no finite subject can guarantee the goodness of its own practice, a subject that could having traversed the is--ought gap and become the infinite reason the series denies. It is the impossibility of that guarantee that constitutes the good as a matter of ethics rather than calculation, and as an unfinished process rather than a possessible state.
